Dos fils conductors rectilinis molt llargs es troben situats paral·lels a l'eix $z$ (perpendiculars al pla $xy$) i separats una distància de 16,0~cm. Pel fil 1 circula un corrent $I_1 = 1,50\ \mathrm{A}$ dirigit cap avall i pel fil 2 circula un corrent $I_2 = 1,50\ \mathrm{A}$ dirigit cap amunt.

\figura[0.6]{fig1.png}

\begin{apartat}{1,25}
Donat un punt $P$ situat al pla $xy$ i equidistant als dos fils (la distància a cada fil és de 10,0~cm), feu un esquema sobre el pla $xy$ del camp magnètic creat per cada fil al punt $P$ i justifiqueu la direcció i el sentit del camp. Quin és el sentit i la direcció del camp magnètic total al punt $P$? Justifiqueu la resposta.
\end{apartat}

\begin{apartat}{1,25}
Calculeu la força per unitat de longitud que fa el fil 1 sobre el fil 2, és a dir, la força que fa el fil 1 sobre un tram d'1,00~m de longitud del fil 2. Indiqueu tant el mòdul com la direcció i el sentit de la força.
\end{apartat}

\dades{%
  El mòdul de camp magnètic creat per un fil infinit per on circula un corrent $I$ a una distància $r$ del fil és $B = \dfrac{\mu_0 I}{2\pi r}$.\\
  $\mu_0 = 4\pi\times 10^{-7}\ \mathrm{T\,m\,A^{-1}}$.}
